\documentclass{article}
\usepackage{amsmath,amssymb}
\usepackage{xurl}
\usepackage[hidelinks]{hyperref}
% Shared commands for notes in docs/paper-gaps/.

\DeclareUnicodeCharacter{2080}{\ifmmode{}_{0}\else\textsubscript{0}\fi}
\DeclareUnicodeCharacter{2081}{\ifmmode{}_{1}\else\textsubscript{1}\fi}
\DeclareUnicodeCharacter{2082}{\ifmmode{}_{2}\else\textsubscript{2}\fi}
\DeclareUnicodeCharacter{2099}{\ifmmode{}_{n}\else\textsubscript{n}\fi}

\newcommand{\Lean}{\textsc{Lean}}
\ExplSyntaxOn
\NewDocumentCommand{\leanid}{m}{
  \ifmmode
    \textsf{\detokenize{#1}}
  \else
    \group_begin:
    \urlstyle{sf}
    \tl_set:Nn \l_tmpa_tl {#1}
    \tl_replace_all:Nnn \l_tmpa_tl {\_} {_}
    \exp_args:NV \path \l_tmpa_tl
    \group_end:
  \fi
}
\ExplSyntaxOff
\providecommand{\tr}{\operatorname{tr}}
\newcommand{\tnleanIssueBase}{https://github.com/LionSR/TNLean/issues/}
\newcommand{\tnleanPrBase}{https://github.com/LionSR/TNLean/pull/}
\newcommand{\tnleanDocBase}{https://sirui-lu.com/TNLean/docs/}
\newcommand{\ghissue}[1]{\href{\tnleanIssueBase#1}{GitHub issue \##1}}
\newcommand{\ghpr}[1]{\href{\tnleanPrBase#1}{PR \##1}}
\newcommand{\doclink}[1]{\href{\tnleanDocBase#1.html}{\path{#1}}}

% Dirac notation and the identity operator, as in blueprint/src/macros/common.tex.
% \providecommand keeps note-local definitions authoritative.
\usepackage{dsfont}
\providecommand{\ket}[1]{|#1\rangle}
\providecommand{\bra}[1]{\langle #1|}
\providecommand{\braket}[2]{\langle #1 | #2 \rangle}
\providecommand{\ketbra}[2]{|#1\rangle\!\langle #2|}
\providecommand{\Id}{\mathds{1}}
\providecommand{\one}{\mathds{1}}
\providecommand{\C}{\mathbb{C}}
\providecommand{\MN}[1]{M_{#1}(\C)}
\providecommand{\MD}{M_D(\C)}

% Verdict marker: \gapnote{<kind>}{<status>}. Kind is the policy
% classification (clarification, local-correction, scope-restriction,
% unfaithful, false-source, open-gap); status is open, wip (elimination
% underway), resolved, or historical. Machine-read by the site index;
% prints a verdict line.
\newcommand{\gapnote}[2]{\par\noindent\textbf{Verdict:} #1 (#2).\par}

\title{The Full-Input Pseudoinverse Obstruction and Its Support Correction}
\author{}
\date{2026-10-04}
\begin{document}
\maketitle
\gapnote{false-source}{resolved}

\section{Source and the full-input obstruction}
For a blocked matrix product map $B=VP$,
arXiv:2307.01696, page~3, footnotes~3 and~4 to equations~(13)--(15),
allows $B$ to be non-injective, interprets the inverse of $P$ as its
pseudoinverse, and permits the input to be placed at any chosen tensor,
with sequential decompositions from both boundaries toward that site.

If $B$ is injective, then $V^\dagger V=\Id_{D^2}$ and the inward sweeps
produce a central isometry on the full virtual-pair input.
For non-injective $B$, its polar support $\Pi$ instead satisfies
$V^\dagger V=\Pi$. Thus a central map on all $D^2$ input coordinates
cannot remain an isometry when $\Pi\ne\Id_{D^2}$.
The assertion that the full-input isometry identity in equation~(15)
remains valid under footnote~3 is therefore false as printed. Restricting
the domain to the support corrects the assertion; it does not prove that
full-input identity.

\section{A nonzero rank-deficient witness}
Take $d=D=2$ and the two nonzero matrices
\begin{align}
    A^0 & =\begin{pmatrix}1&0\\0&0\end{pmatrix}, &
    A^1 & =\begin{pmatrix}0&0\\0&1\end{pmatrix}.
    \label{eq:mswc24_ghz_tensor}
\end{align}
On a ring these give the nonzero vector $\ket{0}^{\otimes N}+\ket{1}^{\otimes N}$.
Block two sites, and order both the physical and virtual pairs as
$00,01,10,11$. Since $B_{ij,ab}=(A^iA^j)_{ab}$, the blocked map and its
canonical polar factors are
\begin{align}
    B & =P=P^+=V=\Pi=\operatorname{diag}(1,0,0,1).
    \label{eq:mswc24_ghz_polar}
\end{align}
Thus $d^2=D^2=4$, both tensor components are nonzero, and
$V^\dagger V=\Pi\ne\Id_4$. The pseudoinverse construction already fails
on these nondegenerate data. One may separately complete the partial
isometry on its kernel, but that completion is not $BP^+$ and is not the
construction asserted in equation~(13).

\section{A nonzero normal counterexample}
Take one physical site with $d=4$ and bond dimension $D=2$, so the
capacity condition $d^q\ge D^2$ in footnote~1 holds with equality. With
$a=1/\sqrt2$, let
\begin{align}
    B^0 & =a\Id, & B^1 & =aE_{01}, & B^2 & =aE_{10}, & B^3 & =0.
    \label{eq:mswc24_pseudoinverse_witness}
\end{align}
This tensor is nonzero, unital and trace preserving. Its length-two
products include $E_{00}/2,E_{01}/2,E_{10}/2,E_{11}/2$, so it is normal.
The one-site physical matrix, in virtual-pair order $00,01,10,11$, is
\begin{align}
    B      & =\begin{pmatrix}a&0&0&a\\0&a&0&0\\0&0&a&0\\0&0&0&0\end{pmatrix}, &
    V=BP^+ & =\begin{pmatrix}a&0&0&a\\0&1&0&0\\0&0&1&0\\0&0&0&0\end{pmatrix}.
    \label{eq:mswc24_pseudoinverse_witness_matrix}
\end{align}
For the unit vector $w=(1,0,0,-1)^{\mathsf T}/\sqrt2$, one has $Vw=0$;
indeed $V^\dagger V=\Id-|w\rangle\langle w|\ne\Id$.
Thus neither a zero tensor nor an insufficient output dimension accounts
for the obstruction. The formal witness is
\leanid{MPSPreparation.mixedPolarFullInputCounterexample_spec}.
Its proof uses the zero final physical row and the actual polar
right-factor identity $V=BG$, without assuming a formula for a shadow
polar map.

An isometric extension on the unused input direction can exist here,
but it differs from $BP^+$. The counterexample concerns that printed
pseudoinverse identification, not the existence of an extension or the
injective-block preparation theorem.

\section{The support construction}
Choose orthonormal coordinates for the range of $\Pi$: an isometry
$J\colon\C^s\to\C^{D^2}$ with $JJ^\dagger=\Pi$.
The bound $s\le D^2$ follows from injectivity of $J$. Then
\begin{align}
    (VJ)^\dagger(VJ) & =J^\dagger\Pi J=\Id_s.
    \label{eq:mswc24_support_isometry}
\end{align}
The existing polar right-factor theorem gives $V=BG$ for a virtual-pair
matrix $G$. Attach $GJ$ at the chosen central site and run the two
independent inward sweeps. Their outer evaluations are isometries
$E_L,E_R$, and their remainders give a central map $C$ such that
\begin{align}
    VJ & =(E_L\otimes\Id_d\otimes E_R)C, & C^\dagger C & =\Id_s.
    \label{eq:mswc24_mixed_support}
\end{align}
Both outer bonds have dimension one, and every side bond has dimension
at most $D^2$. The input may be attached at either endpoint or any
interior site; the argument also allows zero-dimensional support.

The original polar partial isometry is recovered exactly:
\begin{align}
    (VJ)J^\dagger & =V\Pi=V.
    \label{eq:mswc24_support_reconstruction}
\end{align}
Consequently the blocked tensor is $B=(VJ)(J^\dagger P)$. The central
operation is an actual isometry on the compressed input, and no physical
amplitude is discarded by choosing those input coordinates.

\section{Conclusion}
The mixed factorization now covers every blocked tensor on its actual
polar support, with no injectivity assumption. The full-input theorems
remain useful specializations for injective blocks. The same mixed
factorization predicate handles arbitrary finite input dimension, so
no second isometry notion or support representation is introduced.
The corrected support-domain theorem is proved, so the local repair is
resolved. The false-source record and the witness in
(\ref{eq:mswc24_ghz_polar}) remain: the printed full-input pseudoinverse
claim is not established, and the unreduced polar map is only a partial
isometry.\footnote{Formal
    declarations: \leanid{MPSPreparation.exists_mixed_sequential_polar_support_of_split}
    and \leanid{MPSPreparation.exists_mixed_sequential_polar_support} in
    \doclink{TNLean/MPS/Preparation/MixedSequentialSupport}.}
\end{document}
